\section{实证结果}

\begin{frame}{基准回归结果：绿色金融对碳排放强度有显著抑制作用}
  \begin{table}
    \centering
    \small
    \begin{tabular}{@{}lccc@{}}
      \toprule
      模型列 & (1) & (2) & (3) \\
      \midrule
      绿色金融发展指数 & 负 & 负 & \textbf{-2.363*} \\
      经济发展水平控制 & 否 & 是 & 是 \\
      全部控制变量 & 否 & 否 & 是 \\
      固定效应 & 双向 & 双向 & 双向 \\
      \bottomrule
    \end{tabular}
  \end{table}

  \begin{alertblock}{核心发现}
    在纳入全部控制变量后，绿色金融发展指数系数为 \hl{-2.363}，`p=0.057`，
    表明绿色金融发展水平每提高 1\%，碳排放强度平均下降约 \hl{2.36\%}。
  \end{alertblock}

  \begin{itemize}
    \item 人均 GDP 系数显著为负，说明经济发展水平提升有助于降低碳排放强度。
    \item 第二产业占比与煤炭消费占比方向符合预期，但在省级层面统计上不够稳定。
  \end{itemize}
\end{frame}

\begin{frame}{内生性与稳健性：核心结论保持稳定}
  \begin{columns}[T,onlytextwidth]
    \column{0.48\textwidth}
    \begin{block}{工具变量检验}
      \begin{itemize}
        \item 工具变量：滞后一期绿色金融发展指数
        \item 第一阶段 F 统计量约为 \hlb{$2.2\times10^4$}
        \item 第二阶段系数：\hl{-2.307***}
      \end{itemize}
    \end{block}

    \column{0.48\textwidth}
    \begin{exampleblock}{稳健性检验}
      \begin{itemize}
        \item 滞后一期解释变量：\hl{-2.242*}
        \item 替换为碳排放总量：\hl{-1.035**}
        \item 剔除直辖市样本：\hl{-3.668***}
      \end{itemize}
    \end{exampleblock}
  \end{columns}

  \begin{alertblock}{结论}
    无论从内生性处理还是多种稳健性检验看，绿色金融对碳排放的抑制效应都没有消失。
  \end{alertblock}
\end{frame}
